
\subsection{\texorpdfstring{\textbf{
      Реализация непрерывной сети Хопфилда в задаче LMAPF}}
  {Реализация непрерывной сети Хопфилда в задаче LMAPF}}

Самым сложным этапом реализации был именно алгоритм минимизации,
поскольку это задача многих переменных -- 20 при предсказывании 10
шагов, 10 на движение, 10 на вращение. При том, очень нестабильный
функционал стоимости, так как управление на нулевом шаге сбивает
стоимости на всех последующих шагах. (Рис. \ref{media:image_lmapf_1})

В исследовании для решения задачи минимизации с ограничениями была
приведена библиотека IPOPT - Interior Point OPTimizer, но увы, она
написана на C++ и на текущий момент нет связки c Rust. Но была другая
библиотека argmin, которая позволяла решать задачу, но там плохая
документация и полное отсутствие ограничений в задаче минимизации. В
итоге была выбрана библиотека OpEn - Optimisation Engine.
\cite{optimization_engine}

Она позволила решать задачу, но начальное использование центральных
разностей в качестве градиента функционала стоимости приводило к очень
долгим вычислениям. 300 миллисекунд на 30 шагов алгоритма, что
совершенно неприемлемо.

Хотелось вычислить градиент аналитически, но безуспешно. Было решено
задать две компоненты -- D и G в качестве ограничений, что следовало
сделать изначально, так как агент ни при каких обстоятельствах не должен
врезаться в других агентов или стены. (Рис. \ref{media:image_lmapf})


\begin{figure}[h]
  \centering
  \includegraphics[width=1.0\textwidth]{image_lmapf_1.png}
  \caption{}
  \label{media:image_lmapf_1}
\end{figure}


\begin{figure}[h]
  \centering
  \includegraphics[width=0.8\textwidth]{image_lmapf.png}
  \caption{}
  \label{media:image_lmapf}
\end{figure}

Так, идея не была доведена до конца.

\clearpage